% mazzi: 04
\chapter{Validità degli strumenti di misura}

Dopo i presupposti scientifici della valutazione funzionale, un altro aspetto fondamentale è la
\textbf{validità degli strumenti di misura} impiegati nella valutazione. La vera difficoltà non sta
tanto nel somministrare il test, quanto nel \textbf{ridurre il più possibile gli errori} --- di
metodo, di esecuzione e di misurazione --- che in alcune situazioni sono quasi fisiologici.

% ---------------------------------------------------------------------------
\section{Accuratezza e precisione di una misura}

\subsection{Accuratezza}

\textbf{Accuratezza}: quantifica la \textbf{distorsione} in una misura causata da errori di
misurazione.

\begin{itemize}
  \item è quantificata come la \textbf{differenza tra un valore vero e un valore atteso /
    osservato};
  \item una misurazione effettuata con alta precisione avrà sempre un \textbf{piccolo errore}, cioè
    una piccola distorsione dal valore vero;
  \item viene normalmente indicata come l'\textbf{errore massimo} che può esistere in una
    misurazione, ma in realtà ciò che si quantifica è l'\textbf{imprecisione}, misurando la
    \textbf{deviazione tra le misure};
  \item non riguarda quindi la \textbf{media} tra le misure, bensì la \textbf{deviazione standard}
    delle misure, cioè la \textbf{variabilità} tra una misura e l'altra.
\end{itemize}

\subsection{Precisione}

\textbf{Precisione}: è la \textbf{differenza tra un valore osservato e un valore medio atteso}.

\begin{itemize}
  \item se un sistema misura qualcosa che dovrebbe restare invariata, una \textbf{variazione nei
    valori misurati} indica una \textbf{mancanza di precisione};
  \item in genere viene quantificata utilizzando la \textbf{deviazione standard} delle misure
    \textbf{ripetute} della stessa quantità;
  \item precisione e accuratezza sono a volte confuse, ma sono \textbf{quantitativi distinti}.
\end{itemize}

\rubrica{Esempio: il podobarografo a sensore singolo}
Misurando il \textbf{picco di pressione sotto i piedi} durante il passo con un podobarografo
costituito da un \textbf{solo sensore}:
\begin{itemize}
  \item il sistema sarà \textbf{molto preciso}, ma solo \textbf{in quel punto};
  \item sarà invece \textbf{lontano dalla realtà}, perché si perde tutta l'informazione sulla
    \textbf{distribuzione} della pressione sotto il piede nel suo complesso.
\end{itemize}

\subsection{Combinazioni di accuratezza e precisione}

La figura delle slide illustra le tre permutazioni possibili in funzione del \textbf{bersaglio},
che rappresenta il \textbf{valore vero}.

\begin{enumerate}
  \item[(a)] \textbf{alta precisione e alta accuratezza}: la misurazione centra l'obiettivo in
    maniera completa, con entrambe le caratteristiche;
  \item[(b)] \textbf{alta precisione e bassa accuratezza}: la variabilità delle misurazioni è
    piccolissima, ma l'insieme delle misure è \textbf{lontano dal valore vero};
  \item[(c)] \textbf{bassa precisione e scarsa accuratezza}: non si verifica nessuna delle due
    condizioni --- la variabilità tra le misure è enorme e l'accuratezza è scarsa.
\end{enumerate}

\rubrica{Nota sulla slide}
Nella slide la voce (c) riporta ``bassa precisione e scarsa precisione'': il docente segnala
l'\textbf{errore} e corregge, precisando che il secondo termine sta per \textbf{accuratezza}.

% ---------------------------------------------------------------------------
\section{Grandezze misurabili}

Tutte concorrono, a vario titolo, alla valutazione funzionale delle capacità motorie; le prime due
famiglie derivano più o meno direttamente dalla \textbf{fisica}.

\begin{itemize}
  \item \textbf{grandezze cinematiche}: tempo, distanza, velocità, accelerazione;
  \item \textbf{grandezze dinamiche}: forza, coppia, potenza, lavoro, impulso, quantità di moto;
  \item \textbf{grandezze fisiologiche}: FC, $VO_2max$, lattato ematico, EMGs, temperatura.
\end{itemize}

% ---------------------------------------------------------------------------
\section{Il processo di misurazione strumentale}

Gli strumenti sfruttano la \textbf{variazione di tensione} che si sviluppa in seguito alla
prestazione misurata:

\begin{center}
grandezza misurata $\rightarrow$ trasduttore $\rightarrow$ tensione
\end{center}

\rubrica{Trasduttori}
Potenziometri, elettrogoniometri, encoder, celle di carico, pedana dinamometrica, accelerometri.

\begin{itemize}
  \item p.\,es.\ le \textbf{celle di carico}, sottoposte a forza, subiscono una \textbf{piccola
    deformazione} che crea una \textbf{variazione di tensione}; attraverso il trasduttore questa
    fornisce la grandezza misurata, secondo il tipo di scala e le unità di misura utilizzate;
  \item i valori misurati dalla strumentazione \textbf{contengono normalmente errori}, anche nelle
    strumentazioni più precise.
\end{itemize}

% ---------------------------------------------------------------------------
\section{Proprietà dei trasduttori}

Le seguenti proprietà dei trasduttori possono \textbf{influenzare la validità della misura}.

\subsection{Sensibilità}

\textbf{Sensibilità}: riguarda la \textbf{relazione tra il segnale di uscita e quello di entrata};
è data dal \textbf{rapporto tra la variazione della grandezza in ingresso e la variazione della
grandezza in uscita}.

\begin{itemize}
  \item un trasduttore è \textbf{molto sensibile} quando a una \textbf{piccola variazione
    dell'ingresso} corrisponde una \textbf{grande variazione dell'uscita};
  \item è \textbf{poco sensibile} quando a un'\textbf{elevata variazione del segnale di ingresso}
    l'uscita risponde con una \textbf{piccola variazione};
  \item p.\,es.\ una piattaforma di forza con sensibilità di \textbf{10\,N/v}: ogni 10\,N di forza
    generano una tensione di 1\,V (quindi 20\,N $=$ 2\,V). In questo caso la sensibilità è
    \textbf{ridotta}.
\end{itemize}

\subsection{Errore di non linearità}

\textbf{Errore di non linearità}: teoricamente, quando la grandezza misurata aumenta, anche
l'output di tensione deve \textbf{crescere linearmente}; lo \textbf{scostamento dalla linearità}
rappresenta l'errore dello strumento.

\begin{itemize}
  \item p.\,es.\ la \textbf{cella di carico} sottoposta a una forza esterna si deforma in modo
    \textbf{proporzionale} --- quindi lineare --- all'aumentare della forza applicata; è questa
    relazione a permettere la \textbf{trasformazione dei volt} (la variazione di tensione subita
    dallo \textit{strain gauge}) nella forza esterna applicata;
  \item l'errore è di solito fornito \textbf{direttamente dal costruttore} ed è espresso come
    \textbf{spostamento massimo dalla retta ideale};
  \item in condizioni di linearità lo strumento \textbf{non presenta questo errore}.
\end{itemize}

\subsection{Isteresi}

\textbf{Isteresi}: si verifica quando, per \textbf{valori crescenti} della grandezza in ingresso, la
caratteristica assume un determinato andamento, mentre per \textbf{valori decrescenti} della
medesima grandezza ne disegna uno \textbf{differente}.

\begin{itemize}
  \item un buon trasduttore è caratterizzato da un'isteresi \textbf{praticamente nulla}: la
    variazione all'\textbf{aumentare} della grandezza fisica misurata deve essere \textbf{uguale} a
    quella al suo \textbf{diminuire};
  \item molte strumentazioni presentano quasi sempre una \textbf{leggera isteresi}.
\end{itemize}

\subsection{Cross-talk}

\textbf{Cross-talk}: si verifica solitamente nei trasduttori \textbf{triassiali} (accelerometri,
pedane di forza) quando, esercitando una forza su un \textbf{solo asse} (p.\,es.\ l'asse $z$), si
osserva una \textbf{variazione di segnale anche nelle altre direzioni}.

\begin{itemize}
  \item deve essere \textbf{ridotto al minimo}: altrimenti si crea \textbf{interferenza} e quindi
    una \textbf{variabilità enorme} della strumentazione.
\end{itemize}

\subsection{Risoluzione}

\textbf{Risoluzione}: per ogni strumento di misura vi è un \textbf{limite alla grandezza misurata}
che produce un cambiamento nell'uscita dello strumento; la risoluzione riflette perciò la
\textbf{``finezza''} con cui una misurazione può essere effettuata.

\begin{itemize}
  \item spesso è espressa come \textbf{valore} o come \textbf{percentuale del range di misura
    completo} dello strumento;
  \item con strumenti elettronici, riportare misure con \textbf{molti decimali non indica
    necessariamente un'alta risoluzione}: una grande componente di quei valori può riflettere il
    \textbf{rumore del sistema}.
\end{itemize}

\subsection{Temperatura}

\textbf{Temperatura}: anche la temperatura dell'\textbf{ambiente} in cui si eseguono le misurazioni
può determinare \textbf{variazioni del segnale elettrico in uscita}.

\begin{itemize}
  \item per confrontare misurazioni effettuate in \textbf{tempi diversi} è necessario che siano
    eseguite nelle \textbf{stesse condizioni ambientali};
  \item il motivo è che i \textbf{metalli} che costituiscono gran parte della strumentazione
    elettronica sono sensibili alla temperatura e \textbf{variano le proprie condizioni fisiche} al
    variare di essa;
  \item nella pratica un test svolto sul campo d'estate o a inizio primavera è completamente diverso
    da uno svolto in inverno o a inizio preparazione fisica: di questo aspetto va tenuto conto.
\end{itemize}

% ---------------------------------------------------------------------------
\section{La calibrazione}

\textbf{Calibrazione}: procedura che \textbf{non è diretta a eliminare gli errori}, ma dovrebbe
aiutare a \textbf{tenerli al minimo}; sfrutta un \textbf{modello} che relaziona un \textbf{ingresso
noto} allo strumento di misurazione con un \textbf{output desiderato}.

\begin{itemize}
  \item p.\,es.\ la calibrazione di un \textbf{encoder lineare} consiste nel fornire la misura di
    una \textbf{distanza certa e assoluta}, che la strumentazione metterà in relazione al proprio
    sistema di misura;
  \item questi modelli \textbf{non sono rappresentazioni perfette} del sistema: anche durante
    l'acquisizione dei dati gli errori possono provenire da \textbf{fonti diverse};
  \item gli errori nascono perché le \textbf{condizioni cambiano} rispetto a quando la calibrazione
    è stata eseguita --- p.\,es.\ la \textbf{temperatura dei sensori}; per alcune strumentazioni la
    calibrazione va perciò rifatta \textbf{sistematicamente il giorno stesso del test};
  \item per determinare molti parametri biomeccanici è necessario \textbf{combinare parametri e
    variabili di fonti diverse}: in questo caso gli errori sorgono attraverso la \textbf{propagazione
    degli errori} nella loro combinazione (p.\,es.\ per calcolare la \textbf{quantità di moto} devono
    essere note massa e velocità dell'oggetto);
  \item la buona pratica sperimentale è \textbf{minimizzare gli errori alla fonte}.
\end{itemize}

% ---------------------------------------------------------------------------
\section{Evoluzione e applicazione pratica dell'encoder}

\textbf{Encoder}: \textbf{sensore digitale di posizione} di tipo \textbf{rotativo}.

\begin{itemize}
  \item fornisce una rappresentazione \textbf{temporale}, \textbf{numerica} e \textbf{diretta} di
    \textbf{rotazioni angolari}, cioè la variazione della distanza nel tempo;
  \item attraverso vari meccanismi fornisce anche una \textbf{misura diretta di spostamenti
    rettilinei}.
\end{itemize}

\subsection{Struttura e funzionamento}

È formato da un \textbf{disco} (metallo, ceramica, ecc.) che ruota tramite un'\textbf{asta}
(albero), da un \textbf{LED} (laser) e da un \textbf{rilevatore}.

\begin{itemize}
  \item si tratta di strumentazione \textbf{miniaturizzata}: il disco è molto piccolo e presenta
    lungo la circonferenza dei \textbf{fori} a distanza talmente ridotta tra loro da poterli
    considerare \textbf{lineari};
  \item l'\textbf{emettitore ottico} colpisce il rilevatore in maniera \textbf{alternata}, al
    passaggio dei punti vuoti e di quelli pieni;
  \item può lavorare in due modalità:
    \begin{itemize}
      \item \textbf{in trasparenza}: i segmenti del disco sono alternativamente \textbf{opachi e
        trasparenti};
      \item \textbf{in riflessione}: i segmenti del disco sono alternativamente \textbf{opachi e
        riflettenti};
    \end{itemize}
  \item il \textbf{conteggio del numero di segnali} rilevati dà il numero di passaggi effettuati dal
    disco: \textbf{calibrando} questi passaggi con una \textbf{misura certa} (un metro) si stabilisce
    la \textbf{distanza percorsa} da qualsiasi oggetto a cui l'encoder sia attaccato.
\end{itemize}

\subsection{Applicazione: encoder su macchina isotonica}

Applicando l'encoder al \textbf{pacco pesi} di una \textit{leg extension}, dai tracciati di
\textbf{posizione} e \textbf{velocità} in funzione del tempo si legge l'intero gesto:

\begin{itemize}
  \item \textbf{inizio del movimento}: posizione registrata $=$ 0 e velocità $=$ 0;
  \item \textbf{metà del movimento} (in questo caso $45^\circ$): lo spostamento registrato
    dall'encoder, essendo lineare, corrisponde al \textbf{picco di velocità};
  \item \textbf{massima escursione articolare} ($180^\circ$, gamba distesa): la velocità torna a
    \textbf{zero};
  \item \textbf{ritorno alla fase iniziale}: posizione $=$ 0 e velocità $=$ 0.
\end{itemize}

\rubrica{Dalla posizione ai parametri dinamici}
Avendo la posizione in funzione del tempo si ricava, in cascata:
\begin{itemize}
  \item la \textbf{velocità} (prima misura calcolata);
  \item dalla velocità in funzione del tempo, l'\textbf{accelerazione};
  \item con una \textbf{massa nota}, la \textbf{forza} ($F = m \cdot a$);
  \item da forza e velocità, la \textbf{potenza} ($P = F \cdot v$).
\end{itemize}

Con un sistema semplice si ottengono così, per qualsiasi movimento effettuato su una macchina
isotonica o su un castello, il valore di posizione e \textbf{tutti i parametri cinematici e
dinamici}.

% ---------------------------------------------------------------------------
\section{Grandezza misurata e grandezze calcolate}

Per ridurre l'errore su qualsiasi strumento di misura occorre \textbf{conoscere qual è la grandezza
fisica realmente misurata}: solo così si capisce quali \textbf{calcoli} effettuerà poi il sistema o
il software. Si distinguono perciò una \textbf{grandezza fisica misurata} e le \textbf{grandezze
fisiche calcolate}.

\begin{itemize}
  \item \textbf{encoder}: la grandezza misurata è lo \textbf{spostamento}; si procede per
    \textbf{derivazione} --- $spostamento \rightarrow velocità \rightarrow accelerazione$ (derivata
    prima e derivata seconda);
  \item \textbf{accelerometro}: la grandezza misurata è l'\textbf{accelerazione}; si procede in
    senso inverso, per \textbf{integrazione} --- $accelerazione \rightarrow velocità \rightarrow
    spostamento$ (integrazione prima e seconda).
\end{itemize}

\rubrica{Propagazione degli errori}
\begin{itemize}
  \item l'errore sulla grandezza misurata \textbf{si propaga} alle grandezze calcolate, in funzione
    della \textbf{formula}, del \textbf{livello dell'equazione} e del \textbf{grado di potenza
    dell'algoritmo} utilizzato nelle operazioni successive;
  \item nell'encoder un errore minimo sullo spostamento si propaga sulla velocità e, con
    un'equazione di \textbf{secondo grado}, sull'accelerazione;
  \item questo può rendere le \textbf{misure calcolate scarsamente valide};
  \item da qui l'obiettivo fondamentale di tutte le procedure di valutazione: \textbf{ridurre il più
    possibile l'errore della misurazione} alla fonte.
\end{itemize}
